% Extended Abstract

\chapter{Extended Abstract} % Main chapter title

%----------------------------------------------------------------------------------------
\section{Titulo Provisório}
	Sugestões de resolução de incidentes através do uso de LLM’s

\section{Introdução e Enquadramento do Problema}
	A Critical Techworks (CTW) é uma empresa criada através de uma parceria entre a Critical Software e o grupo BMW. Desde 2018 que se dedica ao **desenvolvimento de software** para os vários produtos e serviços digitais da BMW. Tendo em conta a quantidade de sistemas ativos e a sua importância, a forma como os incidentes são geridos torna-se um tema essencial para garantir a estabilidade dos serviços.
	
	Um incidente é, na prática, qualquer evento que cause uma interrupção ou uma degradação no funcionamento de um serviço. Quando isso acontece, é criado um *ticket* com a descrição do problema, a prioridade e outra informação útil para a equipa responsável. Estes incidentes podem resultar de alertas automáticos, baseados em métricas configuradas pelas próprias equipas, ou ser reportados manualmente por utilizadores ou equipas de testes. Em períodos de grande volume de trabalho, ou quando o engenheiro está de prevenção durante a madrugada, é comum surgirem dificuldades na interpretação correta do problema, o que pode levar a erros e aumentar o tempo necessário para encontrar a solução.
	
	O principal desafio identificado é precisamente esta dificuldade em compreender rapidamente o que está a acontecer, sobretudo quando o *ticket* tem pouca informação, é pouco claro ou chega num momento menos favorável. Isto afeta a eficiência das equipas, aumenta o tempo médio de mitigação e pode ter impactos relevantes, tanto a nível financeiro como na experiência de quem utiliza os serviços. Os principais envolvidos neste processo são as equipas de desenvolvimento e operação, os engenheiros de prevenção, os utilizadores internos e externos e a própria organização, que depende da fiabilidade dos seus serviços. O objetivo deste estudo passa por encontrar formas de melhorar a análise dos incidentes e acelerar a sua resolução, reduzindo o impacto que estes problemas podem ter no dia-a-dia da empresa.

\section{Pergunta(s) de Investigação e Objetivos}
	Com base nas dificuldades identificadas pelas equipas \textit{DevOps} e \textit{CorpOps} na análise e resolução de incidentes, surgiu a hipótese de recorrer a ferramentas de Inteligência Artificial, como técnicas de Processamento de Linguagem Natural (\textit{PLN}), para apoiar a compreensão e tratamento destes registos. Desta forma, foi formulada a seguinte questão: “Quão eficazes são os Large Language Models na melhoria da interpretação e análise de registos de incidentes de TI?”.
	
	O objetivo passa por perceber que soluções existem atualmente, de que forma têm sido aplicadas e se há evidências da sua eficácia. Podendo, assim, perceber se existe viabilidade para o desenvolvimento de uma solução que apoie as equipas durante as operações dos seus serviços. Com isso, as equipas poderão conseguir reduzir o Tempo Médio de Resposta (\textit{TMR}) aos incidentes, reduzindo o tempo de indisponibilidade das aplicações, e podendo concentrar maior tempo para o desenvolvimento de outras funcionalidades.

\section{Estado da Arte}
	Optou-se por utilizar um método de pesquisa com base num estudo sistemático de mapeamento (Systematic Mapping Study ou SMS), a fim de, realizar a revisão de literatura. Primeiramente, foi definida uma pergunta de investigação, para estudar soluções já existentes no mercado, e, de seguida, com base na \textit{RQ} (\textit{Research Question}) definida e com o suporte da metodologia \textit{PICOCS}, foram identificados os elementos necessários para estruturar a pesquisa. Com a questão já enquadrada, prosseguiu-se para definição das palavras chave e dos limites de inclusão e exclusão. Por fim, para obter os artigos, foram utilizadas bases de dados de artigos científicos reconhecidas, tais como, a ACM Digital Library e a IEEE Xplore.

	\subsection{Questão}
		Quão eficazes são os Large Language Models na melhoria da interpretação e análise de registos de incidentes de TI?
	
	\subsection{Metodologia PICOC(T)}
		A metodologia PICOC(T) é um pequeno guia que ajuda a organizar e a clarificar uma pergunta de investigação. Com esta metodologia podemos dividir o problema em partes simples, como o público-alvo, a solução em análise, o que será comparado e o contexto, tornando mais fácil de perceber exatamente o que se pretende estudar.
		
		\begin{table}[htbp]
			\centering
			\begin{tabular}{|l|p{3cm}|p{3cm}|p{4cm}|}
				\hline
				\textbf{PICOCT} & \textbf{RQ part} & \textbf{I/E} & \textbf{Sub String} \\
				\hline
				\textbf{P – População} & Registos de incidentes & I — Ferramenta ServiceNow\newline 
				
				E — Estudos sobre gestão de alarmes & "incident management"; "incident analysis"; "incident resolution"; "incident interpretation"; "ticket analysis"; "ticket interpretation"; "ServiceNow" \\
				\hline
				\textbf{I – Intervenção} & Uso de LLMs para análise de incidentes & I — Técnicas conhecidas de LLMs\newline
				
				E — - - - & "Large Language Model"; "LLM"; "Natural Language Processing"; "NLP"; "RAG"; "fine-tuning" \\
				\hline
				\textbf{C – Comparação} & ——— & ——— & ——— \\
				\hline
				\textbf{O – Outcome} & Redução do MTTR, melhoria na priorização e diminuição de erros de analise & I — - - -\newline
				
				E — Estudos não relacionados com incidentes & "incident"; "ticket" \\
				\hline
				\textbf{C – Contexto} & Ambientes DevOps/CorpOps & I — Estudos na área\newline
				
				E — Estudos fora da área & "DevOps"; "CorpOps"; "IT"\\
				\hline
				\textbf{T – Tipo de estudo} & Pesquisa empírica ou com base em estudos & I — - - -\newline
				
				E — Trabalho teórico ou especulativo & Exclude —> "theoretical"; "speculative" \\
				\hline
			\end{tabular}
			\caption{Metodologia PICOCS}
			\label{tab:picocs}
		\end{table}


	\subsection{Palavras Chave Utilizadas}
		\lstset{language=Java}
		\begin{lstlisting}
			("incident management" OR "incident analysis" OR "incident resolution" OR "incident interpretation" OR "ticket analysis" OR "ticket interpretation" OR "ServiceNow")
			AND
			("Large Language Model" OR "LLM" OR "Natural Language Processing" OR "NLP" OR "RAG" OR "fine-tuning")
			AND
			("incident" OR "ticket")
			AND
			("DevOps" OR "CorpOps" OR "IT")
			NOT
			("theoretical" OR "speculative")
		\end{lstlisting}
	
	\subsection{Critérios de Inclusão e Exclusão}
		Para reduzir o “ruído” resultante das bases de dados e melhorar a eficácia da revisão de literatura, foi decidido \textbf{excluir} artigos com mais de 5 anos e \textbf{incluir} apenas os que eram provenientes das conferências mais relacionadas com o tema, como a International Conference on Software Engineering (ICSE), a ACM Joint European Software Engineering Conference and Symposium on the Foundations of Software Engineering (ESEC/FSE) e a conferência Automated Software Engineering (ASE). 
		
	\subsection{PRISMA}
		Para organizar a revisão de literatura, recorreu-se ao método \textbf{\textit{PRISMA}}, que ajuda a tornar o processo de seleção de estudos \textbf{mais claro e transparente}. No total, foram \textbf{identificados 61 artigos} (provenientes do ACM e IEEE Xplore), dos quais \textbf{1 foi removido por ser duplicado}. Seguiu-se a filtragem por tipo de publicação, \textbf{excluindo 26 artigos} de conferências que não se enquadravam no âmbito pretendido. Depois, a leitura dos resumos levou à \textbf{remoção de mais 29 artigos }por falta de relevância. No final, \textbf{permaneceram 5 artigos} que cumprem os critérios definidos e servem de base à análise realizada.
		
	\subsection{Principais trabalhos revistos}
		Dos artigos resultantes da metodologia PRISMA, foram analisados os principais objetivos e conclusões dos respetivos estudos.
		
		\subsubsection{Artigo 1 - Exploring LLM-Based Agents for Root Cause Analysis}
			Esta investigação destaca um problema comum vivido pelos OCE's aquando de um incidente. As pessoas em operações têm de determinar a causa do incidente para o poder resolver da forma mais eficaz, a este processo chama-se Root Cause Analysis (RCA). Desta forma, o estudo em questão apura soluções, com o uso de LLM's, para automatizar este processo de análise gerando sugestões aos engenheiros que permitam uma mitigação dos incidentes mais eficiente. Este estudo apoia-se não só no histórico de instruções como também em ferramentas complementares de registos de logs. No final do ensaio, o autor afirma que os comentários deixados nos incidentes não têm grande impacto na performance dos agentes de LLM, porém o estudo baseou-se apenas em datasets estáticos o que pode enviesar esta análise.
	
		\subsubsection{Artigo 2 - How to mitigate the incident? an effective troubleshooting guide recommendation technique for online service systems}
			Este estudo, expõe os Guias de mitigação de incidentes (ou TSG - Troubleshooting Guides) como um passo importante na gestão dos incidentes por parte dos OCE's. Tem como objetivo investigar o uso e as características dos TSG's e posteriormente o desenvolvimento de uma ferramenta capaz de identificar os melhores guias de mitigação tendo em conta as descrições dos incidentes. Com esta análise foi possível perceber que ~27\% dos incidentes têm guias correspondentes, o que não se aplica ao nosso dataset, e que ~36\% dos guias ocorriam mais do uma vez por ano, evidenciando a existência de incidentes semelhantes ao longo do ano. Com isto, é possível perceber a importância da análise de incidentes anteriores para a análise e mitigação de incidentes futuros. No final, foi possível concluir que em ~80\% dos incidentes registados, a ferramenta foi capaz de sugerir o TSG mais adequado no Top 1.
		
		\subsubsection{Artigo 3 - Identifying linked incidents in large-scale online service systems}
			Neste artigo, foi apresentado um estudo para a criação da ferramenta LiDAR (Linked Incident identification with DAta-driven Representation), com o objetivo de ligar incidentes semelhantes, duplicados ou com a mesma origem. Os estudos e as experiências realizadas permitiram afirmar que a ferramenta desenvolvida foi efetiva e facilmente aplicada pelos OCE's e que facilitou a gestão de incidentes em sistemas mais complexos.
		
		\subsubsection{Artigo 4 - X-Lifecycle Learning for Cloud Incident Management using LLMs}
			Esta análise focou-se em quantificar a importância de melhorar os prompts realizados aos LLM's com dados reais de diferentes etapas do ciclo de vida do software. É algo importante para o estudo a ser realizado aqui mas mais a longo termo.
		
		\subsubsection{Artigo 5 - Xpert: Empowering Incident Management with Query Recommendations via Large Language Models}
			Esta pesquisa, apesar de ainda se concentrar no tema de gestão de incidentes, foca um pouco mais da sua investigação na recomendação e geração de queries (KQL) para ajudar os OCE's na investigação e mitigação dos incidentes. Novamente, algo que seria um trabalho mais futuro para a proposta inicialmente apresentada neste estudo.
		
	\subsection{Lacuna a Resolver (Alterar titulo)}
		Como é afirmado no artigo [3], identificar corretamente incidentes ligados não só ajuda na mitigação de incidente como também na identificação da sua origem. Isto é importante para o estudo a realizar uma vez que uma das premissas é a sugestão dos passos para mitigar os problemas.
		
		Esta investigação vem, assim, juntar um pouco dos artigos apresentados anteriormente de forma a gerar uma solução capaz de analisar incidentes anteriores e sugerir passos para mitigar os incidentes de uma forma eficaz, isto incluindo também a análise de registos dos próprios serviços.

\section{Metodologia Proposta}


\section{Plano de Trabalho}


\section{Considerações Éticas e Profissionais}


\section{Referências Bibliográficas}
